%================================================================
%  PREÁMBULO — REPORTES ESCOM · IPN
%  Motor: XeLaTeX (o LuaLaTeX) + Biber
%  Uso:   \input{preambulo} tras \documentclass[12pt]{article}
%         y define \primertitulo / \segundotitulo para el encabezado.
%================================================================

%----------------------------------------------------------------
% 1. COLORES (xcolor se carga primero para evitar choques)
%----------------------------------------------------------------
\usepackage[table,svgnames]{xcolor}

% Identidad institucional
\definecolor{ipnguinda}{RGB}{111,17,53}
\definecolor{ipnguindaclaro}{RGB}{155,60,95}
\definecolor{ipngray}{RGB}{90,90,90}
\definecolor{ipngrisclaro}{RGB}{225,225,228}
\definecolor{accent}{RGB}{0,90,160}
\definecolor{ipnverde}{RGB}{30,120,70}
\definecolor{ipnambar}{RGB}{190,120,0}

% Código (tema claro)
\definecolor{codebg}{RGB}{248,248,250}
\definecolor{codefg}{RGB}{40,40,40}
\definecolor{codeframe}{RGB}{218,218,224}
\definecolor{codenum}{RGB}{160,160,168}
\definecolor{codekw}{RGB}{0,92,197}
\definecolor{codetype}{RGB}{111,66,193}
\definecolor{codestr}{RGB}{34,134,58}
\definecolor{codecom}{RGB}{120,128,138}
\definecolor{codegray}{rgb}{0.95,0.95,0.95}

% Terminal (tema oscuro)
\definecolor{termbg}{RGB}{28,28,34}
\definecolor{termbar}{RGB}{48,48,56}
\definecolor{termtitle}{RGB}{170,170,180}
\definecolor{termcmd}{RGB}{255,160,70}
\definecolor{termout}{RGB}{125,220,145}
\definecolor{bg-dark}{rgb}{0.2,0.2,0.2}
\definecolor{fg-white}{rgb}{1,1,1}

%----------------------------------------------------------------
% 2. TIPOGRAFÍA
%----------------------------------------------------------------
\usepackage{fontspec}
\defaultfontfeatures{Ligatures=TeX}
\IfFontExistsTF{Times New Roman}
{\setmainfont{Times New Roman}}
{\setmainfont{TeX Gyre Termes}}          % respaldo métrico idéntico
\IfFontExistsTF{Consolas}
{\setmonofont{Consolas}[Scale=MatchLowercase]}
{\setmonofont{DejaVu Sans Mono}[Scale=MatchLowercase]}

\usepackage{microtype}

%----------------------------------------------------------------
% 3. IDIOMA
%----------------------------------------------------------------
\usepackage[spanish, es-tabla, es-noshorthands]{babel}
\usepackage{csquotes}

%----------------------------------------------------------------
% 4. PÁGINA Y ESPACIADO
%----------------------------------------------------------------
\usepackage{geometry}
\geometry{
    letterpaper,
    left=2.5cm,
    right=2.5cm,
    top=2cm,
    bottom=2cm,
    includehead,
    includefoot,
    headheight=15pt,
    headsep=14pt,
    footskip=30pt
}

\usepackage{setspace}
\onehalfspacing
\usepackage[skip=0.6\baselineskip plus 2pt]{parskip}

% Evita líneas viudas/huérfanas
\clubpenalty=10000
\widowpenalty=10000
\displaywidowpenalty=10000
\raggedbottom

%----------------------------------------------------------------
% 5. MATEMÁTICAS Y UNIDADES
%----------------------------------------------------------------
\usepackage{amsmath}
\usepackage{mathtools}
\usepackage{amssymb}
\usepackage{siunitx}
\sisetup{
    output-decimal-marker = {.},
    group-separator       = {\,},
    group-minimum-digits  = 4,
    group-digits          = integer,
    range-phrase          = { a },
    list-final-separator  = { y },
    list-pair-separator   = { y },
    detect-all
}

%----------------------------------------------------------------
% 6. FIGURAS Y TABLAS
%----------------------------------------------------------------
\usepackage{graphicx}
\graphicspath{{./}{Imagenes/}{imagenes/}{img/}{figuras/}}
\usepackage{float}
\usepackage{array}
\usepackage{booktabs}
\usepackage{longtable}
\usepackage{tabularx}
\usepackage{caption}
\usepackage{subcaption}

% Leyendas estilo APA 7: etiqueta en negritas, título en cursiva abajo
\captionsetup{
    format=plain,
    font={small,stretch=1.1},
    labelfont={bf,color=ipnguinda},
    textfont=it,
    labelsep=newline,
    justification=raggedright,
    singlelinecheck=false,
    skip=6pt
}
\captionsetup[table]{position=top}
\captionsetup[sub]{labelsep=space, labelfont={bf}, textfont=normal, justification=centering}

% Tablas más aireadas y con reglas en tono institucional
\renewcommand{\arraystretch}{1.3}
\setlength{\tabcolsep}{7pt}
\setlength{\aboverulesep}{0pt}
\setlength{\belowrulesep}{0pt}
\setlength{\heavyrulewidth}{1pt}
\setlength{\lightrulewidth}{0.5pt}

% Columnas con ancho fijo y alineación: L{3cm}, C{2cm}, R{2cm}
\newcolumntype{L}[1]{>{\raggedright\arraybackslash}p{#1}}
\newcolumntype{C}[1]{>{\centering\arraybackslash}p{#1}}
\newcolumntype{R}[1]{>{\raggedleft\arraybackslash}p{#1}}
% Para tabularx: Y (izquierda) y Z (centrada) flexibles
\newcolumntype{Y}{>{\raggedright\arraybackslash}X}
\newcolumntype{Z}{>{\centering\arraybackslash}X}

% Encabezado de tabla:  \encabezado \enc{Col 1} & \enc{Col 2} \\
\newcommand{\encabezado}{\rowcolor{ipnguinda}}
\newcommand{\enc}[1]{\textcolor{white}{\bfseries #1}}
% Filas zebra: \zebra antes del tabular (usa \rowcolors de xcolor)
\newcommand{\zebra}{\rowcolors{2}{white}{ipnguinda!5}}

%----------------------------------------------------------------
% 7. LISTAS
%----------------------------------------------------------------
\usepackage{enumitem}
\setlist{itemsep=2pt, topsep=4pt, parsep=0pt, leftmargin=*}
\setlist[itemize,1]{label=\textcolor{ipnguinda}{\small$\bullet$}}
\setlist[itemize,2]{label=\textcolor{ipnguindaclaro}{\scriptsize$\blacktriangleright$}}
\setlist[itemize,3]{label=\textcolor{ipngray}{\footnotesize$\circ$}}
\setlist[enumerate,1]{label=\textcolor{ipnguinda}{\bfseries\arabic*.}}
\setlist[enumerate,2]{label=\textcolor{ipnguindaclaro}{\bfseries\alph*)}}
\setlist[enumerate,3]{label=\textcolor{ipngray}{\roman*.}}
\setlist[description]{font=\color{ipnguinda}\bfseries, leftmargin=1.5em, labelindent=0pt, style=unboxed}

%----------------------------------------------------------------
% 8. CÓDIGO: listings + tcolorbox
%----------------------------------------------------------------
\usepackage{listings}
\usepackage{appendix}
\usepackage[most]{tcolorbox}
\usepackage{mdframed}

% --- Caracteres Unicode > U+00FF dentro de listings (XeLaTeX) ---
% listings solo procesa los códigos 0–255; sin esto, los caracteres de
% árboles (├ │ └ ─), flechas, comillas tipográficas, etc. salen desordenados.
\makeatletter
\newcommand\lstunicodechar[1]{%
    \begingroup\lccode`\~="#1\relax
    \lowercase{\endgroup\def~}{\lst@ProcessOther{\char"#1\relax}}%
    \catcode"#1=\active
}
\newcommand\lstunicodelist{%
    251C,2502,2514,2500,250C,2510,2518,252C,2534,253C,2524,% ├│└─┌┐┘┬┴┼┤
    2550,2551,2554,2557,255A,255D,2560,2563,%                ═║╔╗╚╝╠╣
    2190,2191,2192,2193,21D2,%                                ←↑→↓⇒
    2026,201C,201D,2018,2019,2022,2013,2014,%                 …“”‘’•–—
    2264,2265,2260%                                           ≤≥≠
}
\lst@AddToHook{SelectCharTable}{%
    \@for\lstuni@item:=\lstunicodelist\do{\expandafter\lstunicodechar\expandafter{\lstuni@item}}%
}
\makeatother

% --- Soporte UTF-8 (acentos, ñ, ¿¡) -----------------------------
\lstdefinestyle{estilo_utf8}{
    extendedchars=true,
    keepspaces=true,
    columns=fullflexible,
    showstringspaces=false,
    literate=
    {á}{{\'a}}1 {é}{{\'e}}1 {í}{{\'i}}1 {ó}{{\'o}}1 {ú}{{\'u}}1
    {Á}{{\'A}}1 {É}{{\'E}}1 {Í}{{\'I}}1 {Ó}{{\'O}}1 {Ú}{{\'U}}1
    {ñ}{{\~n}}1 {Ñ}{{\~N}}1 {ü}{{\"u}}1 {Ü}{{\"U}}1
    {¿}{{?`}}1  {¡}{{!`}}1  {°}{{\textdegree}}1
}

% --- Estilo base (tema claro) -----------------------------------
\lstdefinestyle{estilo_base}{
    style=estilo_utf8,
    basicstyle=\ttfamily\small\color{codefg}\linespread{1.05}\selectfont,
    keywordstyle=\color{codekw}\bfseries,
    keywordstyle=[2]\color{codetype},
    commentstyle=\color{codecom}\itshape,
    stringstyle=\color{codestr},
    identifierstyle=,
    numbers=left,
    numberstyle=\tiny\ttfamily\color{codenum},
    stepnumber=1,
    numbersep=9pt,
    tabsize=4,
    breaklines=true,
    breakatwhitespace=true,
    postbreak=\mbox{\textcolor{ipnguindaclaro}{$\hookrightarrow$}\space},
    backgroundcolor=\color{codebg},
    frame=single,
    framerule=0.6pt,
    rulecolor=\color{codeframe},
    framesep=6pt,
    xleftmargin=22pt,
    xrightmargin=4pt,
    aboveskip=10pt,
    belowskip=10pt,
    captionpos=b
}

% --- Estilos por lenguaje (mismos nombres de siempre) -----------
\lstdefinestyle{cstyle}{
    style=estilo_base,
    language=C,
    morekeywords=[2]{size_t,ssize_t,pid_t,uint8_t,uint16_t,uint32_t,uint64_t,
        int8_t,int16_t,int32_t,int64_t,FILE,bool,pthread_t,sem_t,key_t,mqd_t},
    morecomment=[l][\color{ipnguindaclaro}]{\#}
}
\lstdefinestyle{estilo_c}{style=cstyle}

\lstdefinestyle{estilo_bash}{
    style=estilo_base,
    language=bash,
    numbers=none,
    xleftmargin=6pt,
    morekeywords={sudo,apt,make,gcc,git,cd,ls,mkdir,rm,cp,mv,chmod,cat,grep,pip,python,python3}
}

\lstdefinestyle{estilo_make}{
    style=estilo_base,
    language=make,
    tabsize=8
}

\lstdefinestyle{estilo_tree}{
    style=estilo_base,
    language={},
    numbers=none,
    xleftmargin=6pt,
    basicstyle=\ttfamily\small\color{codefg}\linespread{1.0}\selectfont
}

% Estilos adicionales frecuentes
\lstdefinestyle{estilo_python}{style=estilo_base, language=Python}
\lstdefinestyle{estilo_sql}{style=estilo_base, language=SQL, morekeywords={REFERENCES,SERIAL,VARCHAR,BOOLEAN,TIMESTAMP}}
\lstdefinelanguage{JavaScript}{
    keywords={break,case,catch,class,const,continue,default,do,else,export,
        extends,finally,for,from,function,if,import,in,instanceof,let,new,
        return,switch,this,throw,try,typeof,var,while,async,await,of,null,
        true,false,undefined},
    sensitive=true,
    comment=[l]{//},
    morecomment=[s]{/*}{*/},
    morestring=[b]', morestring=[b]", morestring=[b]`
}
\lstdefinestyle{estilo_js}{style=estilo_base, language=JavaScript}
\lstdefinelanguage{HTML5}[]{HTML}{
    morekeywords={header,nav,main,section,article,footer,aside,figure,
        template,canvas,video,audio,source,svg,button,input,label},
    sensitive=false
}
\lstdefinestyle{estilo_html}{style=estilo_base, language=HTML5}

\lstset{style=estilo_base}

% --- Estilo "limpio" para listings dentro de tcolorbox ------------
\lstdefinestyle{estilo_terminal}{
    style=estilo_utf8,
    language=bash,
    numbers=none,
    frame=none,
    backgroundcolor={},
    xleftmargin=0pt,
    xrightmargin=0pt,
    aboveskip=0pt,
    belowskip=0pt,
    keywordstyle={},
    commentstyle={},
    stringstyle={},
    identifierstyle={},
    breaklines=true,
    breakatwhitespace=false,
    postbreak=\mbox{\textcolor{termtitle}{$\hookrightarrow$}\space},
    showtabs=false
}

% --- Terminal estilo ventana (con botones) ----------------------
\tcbset{
    ventana/.style={
        enhanced,
        breakable,
        listing only,
        colback=termbg,
        colframe=termbar,
        colbacktitle=termbar,
        coltitle=termtitle,
        fonttitle=\ttfamily\footnotesize,
        halign title=center,
        toptitle=2pt,
        bottomtitle=2pt,
        titlerule=0pt,
        boxrule=0.4pt,
        arc=6pt,
        left=10pt,
        right=10pt,
        top=6pt,
        bottom=6pt,
        before skip=10pt,
        after skip=10pt,
        overlay unbroken and first={
            \foreach \col/\dx in {red!75!orange/10pt, orange!85!yellow/20pt, green!65!black/30pt}
            \fill[\col] ([xshift=\dx]title.west) circle (2.7pt);
        }
    }
}

\newtcblisting{terminal}[1][Terminal]{
    ventana,
    title={#1},
    listing options={
        style=estilo_terminal,
        basicstyle=\ttfamily\small\color{termcmd}
    }
}

\newtcblisting{terminalresponse}[1][Salida]{
    ventana,
    title={#1},
    listing options={
        style=estilo_terminal,
        language={},
        basicstyle=\ttfamily\small\color{termout}
    }
}

% --- Bloque de código con pestaña de archivo --------------------
%  \begin{codigo}[main.c]{estilo_c} ... \end{codigo}
%  \codigoarchivo[main.c]{estilo_c}{src/main.c}
\tcbset{
    hojacodigo/.style={
        enhanced,
        breakable,
        listing only,
        colback=codebg,
        colframe=codeframe,
        boxrule=0.6pt,
        arc=4pt,
        left=24pt,
        right=6pt,
        top=4pt,
        bottom=4pt,
        before skip=10pt,
        after skip=10pt,
        attach boxed title to top left={xshift=10pt, yshift=-\tcboxedtitleheight/2},
        boxed title style={colback=ipnguinda, colframe=ipnguinda, arc=2pt, boxrule=0pt,
            left=5pt, right=5pt, top=1pt, bottom=1pt},
        fonttitle=\ttfamily\footnotesize\bfseries,
        coltitle=white,
        top=\tcboxedtitleheight/2+4pt
    }
}

\newtcblisting{codigo}[2][]{
    hojacodigo,
    title={#1},
    listing options={style=#2, frame=none, backgroundcolor={}, xleftmargin=0pt,
        xrightmargin=0pt, aboveskip=0pt, belowskip=0pt}
}

\newtcbinputlisting{\codigoarchivo}[3][]{
    hojacodigo,
    title={#1},
    listing file={#3},
    listing options={style=#2, frame=none, backgroundcolor={}, xleftmargin=0pt,
        xrightmargin=0pt, aboveskip=0pt, belowskip=0pt}
}

% Código en línea:  \code{printf()}
\newcommand{\code}[1]{%
    \tcbox[on line, colback=codebg, colframe=codeframe, boxrule=0.4pt, arc=2pt,
    boxsep=0pt, left=2pt, right=2pt, top=1.5pt, bottom=1.5pt,
    fontupper=\ttfamily\small]{\detokenize{#1}}%
}

%----------------------------------------------------------------
% 9. CAJAS DE ÉNFASIS
%----------------------------------------------------------------
%  \begin{nota}[Título opcional] ... \end{nota}
\tcbset{
    cajabase/.style={
        enhanced,
        breakable,
        sharp corners,
        boxrule=0pt,
        leftrule=3pt,
        left=10pt,
        right=10pt,
        top=6pt,
        bottom=6pt,
        before skip=10pt,
        after skip=10pt,
        parbox=false
    }
}

\newcommand{\definircaja}[3]{%  entorno, color, título por defecto
    \newtcolorbox{#1}[1][#3]{
        cajabase,
        colframe=#2,
        colback=#2!6,
        fontupper=\setstretch{1.25}\setlength{\parskip}{4pt},
        before upper={{\bfseries\color{#2}##1}\par}
    }%
}

\definircaja{nota}{accent}{Nota}
\definircaja{importante}{ipnguinda}{Importante}
\definircaja{advertencia}{ipnambar}{Advertencia}
\definircaja{exito}{ipnverde}{Resultado}

% Definiciones numeradas: \begin{definicion}{Proceso} ... \end{definicion}
\newtcbtheorem[number within=section]{definicion}{Definición}{
    enhanced,
    breakable,
    colback=white,
    colframe=ipnguinda,
    colbacktitle=ipnguinda,
    fonttitle=\bfseries,
    separator sign={.},
    boxrule=0.6pt,
    arc=3pt,
    before skip=10pt,
    after skip=10pt
}{def}

%----------------------------------------------------------------
% 10. DIAGRAMAS
%----------------------------------------------------------------
\usepackage{tikz}
\usetikzlibrary{positioning, arrows.meta, shapes.geometric, shapes.misc, fit, calc, shadows, babel}

\tikzset{
    caja/.style={draw=ipnguinda, fill=ipnguinda!6, rounded corners=3pt,
        thick, minimum height=8mm, minimum width=24mm, align=center, font=\small},
    decision/.style={diamond, draw=accent, fill=accent!6, thick, aspect=2,
        align=center, font=\small, inner sep=1pt},
    inicio/.style={rounded rectangle, draw=ipngray, fill=ipngrisclaro, thick,
        minimum height=8mm, minimum width=22mm, font=\small},
    flecha/.style={-{Stealth[length=2.5mm]}, thick, color=ipngray}
}

%----------------------------------------------------------------
% 11. TÍTULOS DE SECCIÓN (16 / 14 / 12 pt)
%----------------------------------------------------------------
\usepackage{titlesec}
\setcounter{secnumdepth}{4}
\setcounter{tocdepth}{4}

% Sección: número en recuadro guinda + filete inferior
\titleformat{\section}[hang]
{\fontsize{16}{19}\selectfont\bfseries\color{ipnguinda}}
{\colorbox{ipnguinda}{\makebox[1.2em]{\color{white}\thesection}}}
{0.6em}
{}
[\vspace{-4pt}{\color{ipnguinda!35}\titlerule[0.8pt]}]
\titleformat{name=\section, numberless}[hang]
{\fontsize{16}{19}\selectfont\bfseries\color{ipnguinda}}
{}{0pt}{}
[\vspace{-4pt}{\color{ipnguinda!35}\titlerule[0.8pt]}]
\titlespacing*{\section}{0pt}{18pt}{10pt}

% Subsección
\titleformat{\subsection}[hang]
{\fontsize{14}{17}\selectfont\bfseries}
{\textcolor{ipnguinda}{\thesubsection}}
{0.7em}
{}
\titlespacing*{\subsection}{0pt}{14pt}{6pt}

% Subsubsección
\titleformat{\subsubsection}[hang]
{\fontsize{12}{14.5}\selectfont\bfseries}
{\textcolor{ipnguindaclaro}{\thesubsubsection}}
{0.7em}
{}
\titlespacing*{\subsubsection}{0pt}{12pt}{4pt}

% Cuarto nivel
\titleformat{\paragraph}[hang]
{\normalfont\normalsize\bfseries\itshape}
{\textcolor{ipngray}{\theparagraph}}
{0.7em}
{}
\titlespacing*{\paragraph}{0pt}{10pt}{2pt}

%----------------------------------------------------------------
% 12. ÍNDICE
%----------------------------------------------------------------
\usepackage{titletoc}

\titlecontents{section}[2em]
{\addvspace{8pt}\bfseries}
{\contentslabel[\textcolor{ipnguinda}{\thecontentslabel}]{2em}}
{\hspace*{-2em}}
{\hfill\textcolor{ipnguinda}{\contentspage}}

\titlecontents{subsection}[4.8em]
{\addvspace{1pt}}
{\contentslabel{2.8em}}
{\hspace*{-2.8em}}
{\titlerule*[0.7pc]{\textcolor{ipngray}{.}}\contentspage}

\titlecontents{subsubsection}[8.4em]
{\small}
{\contentslabel{3.6em}}
{\hspace*{-3.6em}}
{\titlerule*[0.7pc]{\textcolor{ipngray}{.}}\contentspage}

\titlecontents{paragraph}[12.8em]
{\small\itshape}
{\contentslabel{4.4em}}
{\hspace*{-4.4em}}
{\titlerule*[0.7pc]{\textcolor{ipngray}{.}}\contentspage}

\titlecontents{figure}[3.2em]
{\addvspace{2pt}}
{\contentslabel{3.2em}}
{\hspace*{-3.2em}}
{\titlerule*[0.7pc]{\textcolor{ipngray}{.}}\contentspage}
\titlecontents{table}[3.2em]
{\addvspace{2pt}}
{\contentslabel{3.2em}}
{\hspace*{-3.2em}}
{\titlerule*[0.7pc]{\textcolor{ipngray}{.}}\contentspage}

%----------------------------------------------------------------
% 13. ENCABEZADO Y PIE DE PÁGINA
%----------------------------------------------------------------
\usepackage{fancyhdr}
\usepackage{lastpage}

\renewcommand{\sectionmark}[1]{\markboth{\thesection.\ #1}{}}
\renewcommand{\subsectionmark}[1]{\markright{\thesubsection.\ #1}}

% Estilo de texto de encabezado/pie
\newcommand{\hdrfont}{\footnotesize\color{ipngray}}

\newcommand{\numpagina}{\colorbox{ipnguinda}{\makebox[1.4em]{\color{white}\bfseries\thepage}}}

\pagestyle{fancy}
\fancyhf{}
\fancyhead[L]{\hdrfont\textcolor{ipnguinda}{\bfseries ESCOM}\,\textperiodcentered\,IPN}
\fancyhead[C]{\hdrfont\ifdefined\segundotitulo\segundotitulo\fi}
\fancyhead[R]{\hdrfont\itshape\ifdefined\primertitulo\primertitulo\fi}
\fancyfoot[L]{\hdrfont\nouppercase{\leftmark}}
\fancyfoot[R]{\hdrfont\numpagina}
\renewcommand{\headrulewidth}{0.6pt}
\renewcommand{\footrulewidth}{0.4pt}
\renewcommand{\headrule}{{\color{ipnguinda}\hrule height \headrulewidth width \headwidth \vskip-\headrulewidth}}
\renewcommand{\footrule}{{\color{ipngrisclaro}\hrule height \footrulewidth width \headwidth \vskip-\footrulewidth}}

% Páginas "plain" (p. ej. \maketitle) con solo el número
\fancypagestyle{plain}{%
    \fancyhf{}%
    \fancyfoot[R]{\hdrfont\numpagina}%
    \renewcommand{\headrulewidth}{0pt}%
    \renewcommand{\footrulewidth}{0pt}%
}

%----------------------------------------------------------------
% 14. NOTAS AL PIE
%----------------------------------------------------------------
\usepackage[hang, flushmargin]{footmisc}
\renewcommand{\footnoterule}{\kern-3pt{\color{ipnguinda!40}\hrule width 0.3\columnwidth height 0.5pt}\kern 2.5pt}

%----------------------------------------------------------------
% 15. BIBLIOGRAFÍA APA 7
%----------------------------------------------------------------
\usepackage[
backend=biber,
style=apa,
language=spanish
]{biblatex}
\IfFileExists{spanish-apa.lbx}{\DeclareLanguageMapping{spanish}{spanish-apa}}{}
\addbibresource{Referencias/referencias.bib}
%\addbibresource{Referencias/references.bib}
%\addbibresource{References/referenciasdieguito.bib}

\setlength{\bibhang}{1.27cm}              % sangría francesa APA
\setlength{\bibitemsep}{0.4\baselineskip}
\renewcommand*{\bibfont}{\normalsize\setstretch{1.15}}

%----------------------------------------------------------------
% 16. HIPERVÍNCULOS (al final)
%----------------------------------------------------------------
\usepackage{xurl}
\usepackage[
colorlinks=true,
linkcolor=ipnguinda,
urlcolor=accent,
citecolor=ipnguinda,
filecolor=accent,
bookmarksnumbered=true,
bookmarksopen=true,
pdfpagemode=UseOutlines,
pdfstartview=FitH
]{hyperref}
\usepackage{bookmark}

\AtBeginDocument{%
    \renewcommand{\lstlistingname}{Código}%
    \renewcommand{\lstlistlistingname}{Índice de códigos}%
    \hypersetup{
        pdfcreator={XeLaTeX},
        pdfsubject={ESCOM -- IPN}
    }%
    \ifdefined\primertitulo\hypersetup{pdftitle={\primertitulo}}\fi
}